\section{A persistent truth removes the optional-major branch}
\label{sec:mandatory-major}

Let $\mathcal A$ be the explicitly stated native theory with numerical
constructors $0,S,I$, equations, the paired binder $D$, and auxiliary
closure $Q$.  Write $\top_*=(0=0)$ and $\bot_*=(0=S0)$.
The native theory proves $\top_*$ and has elimination from $\bot_*$.
Its full-iPRA interpretation is left unchanged.

\subsection{The persistent truth and the compulsory premise}

Retain the inherited numerical decoder and canonical formula decodings.
The empty array and arrays of arity at least four read as $\top_*$.
Unary, binary, and ternary arrays retain their $D$, equation, and $Q$
readings.  As in the inherited representation, scalar padding at formula
positions may be assigned the harmless truth reading.  The proof below
shows that a raw numeric scalar never succeeds as a checking goal.

The initial context is the one-element list containing the empty array.
Thus it contains one proved truth.  Every non-hypothesis packet now
checks its submitted major.  The canonical absent major is the empty
array; its compulsory proof is an immediate hypothesis use.  In a
paired-introduction packet, the case/progress marker is moved to the
initializer field: $u=[]$ requests cases and $u=0$ requests progress
where the opened matrix is a literal equation.  A nonempty initializer
still selects the inherited computation branch.

The point is not merely to omit a guard.  Python identifies numeric
zero with its false sentinel, so an arbitrary unchecked scalar major
could accidentally match the closure branch.  A persistent truth premise
and the invariant below make that scalar case unreachable on every
successful checking trace.  No input type promise or uncounted wrapper
is assumed.

\subsection{The root-shape invariant}

Call an expression \emph{array rooted} when its outermost Python value
is a list.  The distinction concerns the outer constructor only.
It places no restriction on the arities or contents of that list.

\begin{lemma}[Root preservation]
Whenever the explicit-list opener succeeds, an array input has an array
output and a numeric scalar input has a numeric scalar output at formula
root depth.  Other raw JSON values fail its scalar comparison.
The transport assignment preserves the root class of the current goal.
\end{lemma}
\begin{proof}
The array clause returns a list comprehension.  At formula root depth
zero, a numeric scalar is retained or replaced by scalar zero; no numerical
image is inserted at that root.  Nonnumeric nonarrays cannot pass the
scalar ordering comparison.

For a successful transport match, the target goal equals the opening
of its supplied template at one endpoint.  Both endpoint openings have
the same root class as that template.  Replacing the target by the other
opening therefore preserves the root class. A failed match leaves the
goal unchanged. The compact profile's false-valued source fallback
likewise leaves it unchanged.
\end{proof}

\begin{lemma}[Successful-trace invariant]
\label{lem:successful-array-invariant}
Suppose all entries of the incoming context are array rooted.  If a
finite call to the explicit-list checker succeeds, its original goal
is array rooted, and every newly appended hypothesis in that successful
trace is array rooted.  Every compulsory major in the trace is
therefore array rooted.
\end{lemma}
\begin{proof}
Induct on the finite successful checking tree.  If hypothesis closure
succeeds, the opened goal equals an array-rooted context entry.  Root
preservation implies that the original goal is array rooted.

Otherwise the compulsory recursive call on the major succeeds first.
Its incoming context has only array-rooted entries, so the induction
hypothesis implies that the raw major $m$ is a list.  If $m$ is nonempty,
its export $m*3$ is a list.  If it is empty and computation is selected,
the generated hypothesis is an explicit binary list.  If it is empty
and introduction is selected, the hypothesis and updated goal are the
same opened matrix.  A scalar matrix would fail the subsequent length
operation because the empty major cannot bypass that test.  Hence in
a successful introduction this hypothesis is again array rooted.
The original introduction goal must itself have been an array: its
initial opener rejects nonnumeric nonarrays, and a numeric scalar
cannot be starred into the introduction call.

Consequently the last recursive call has an array-rooted appended
hypothesis and an incoming context consisting entirely of arrays.
Its induction hypothesis implies that its goal is array rooted.
In the major and computation branches, root preservation, including
transport preservation, then implies that the original goal was array
rooted.  The introduction case was already checked.  This establishes
both claims simultaneously.
\end{proof}

The lemma applies at the public entry point because the sole initial
hypothesis is the empty array.  Numeric zero and false can never be
accepted majors, even if supplied adversarially in a packet.  The
argument concerns successful finite traces; it is not a probability
claim derived from negative tests.

\subsection{Opening and freshness}

The inherited opening operation traverses arrays at the same depths as
before.  The numerical cases, capture-free insertion, formula padding,
and array interpretation are unchanged.  Before a new floor is allocated,
the retained context and current relevant goal have passed the no-image
check.  They therefore contain no newly admitted relevant numerical
parameter.  If a finite input contains fractional negative names, all
additional names in the new interval are fresh.  The zero-substitution
argument below substitutes zero for each of them, not merely for the
integer endpoint.

\begin{theorem}[Soundness of the mandatory-major presentation]
Every finite successful run of the explicit-list checker has a decoded
conclusion derivable in $\mathcal A$ from its decoded retained
assumptions.  In particular, every closed conclusion accepted from the
initial context is a theorem of $\mathcal A$.
\end{theorem}
\begin{proof}
Apply Lemma~\ref{lem:successful-array-invariant} and induct on the
finite successful checking trace.  Hypothesis closure uses semantic
opening.  Every other packet proves an array-rooted major first.
There are exactly two possibilities.

\paragraph{Nonempty array major.}
The inherited major argument applies.  A match with the binary array
$[s,u]$ is equality transport.  A match with $[t,t,t]$ when $s$ is
false-valued is closure specialization from its zero instance.  A
mismatch retains the old goal. For the compact profile, if its Boolean
expression declines to
replace a goal by a false-valued source representation, the unchanged
goal is still checked, which only strengthens the obligation.

Tripling a unary major yields the canonical three-copy closure of its
matrix, and the continuation is paired elimination with the fresh witness.
Tripling an array with at least two entries yields an array with at least
six entries, hence a proved truth alias.  The inherited opening and
freshness lemmas justify the continuation.  The numeric false sentinel
cannot match this nonempty array major.

\paragraph{Empty array major.}
Its compulsory premise is $\top_*$.  An empty array cannot equal the
binary equality, the nonempty triple closure, or the numeric false
sentinel.  Thus the transport test leaves the original goal unchanged.
The empty export selects computation or introduction.

For introduction, exact unary unpacking handles a genuine $D$ goal.
The witness is opened, the index coordinate retained, and the continuation
checks the successor instance.  When $u=0$, the test
$\operatorname{len}(a)\bmod2=u$ retains the arbitrary-index hypothesis
only for an even-arity opened body.  Among nontrivial formula constructors
this is precisely a literal equation.  Empty and larger even arities mean
truth, where extra progress instances are admissible.  Unary $D$ and
ternary $Q$ bodies receive the zero-instance hypothesis, as required by
cases.  The marker $u=[]$ always selects cases.  Other successfully used
false-valued numeric markers have the same numerical-zero role.  A scalar
opened body fails its length operation.  The undersupplied empty-goal case
already has conclusion $\top_*$; exact unpacking rejects a binary goal
and arrays with more than one argument.  In particular no nontrivial
compound induction rule is admitted.

For computation, the generated hypothesis is a binary equation.  Its
all-zero fresh-name instance is the inherited valid computation instance,
including the certified-zero repair where needed.  The test involving
$u$ can select an open fresh instance instead of its zero instance.
The open equation need not be universally true.  Since the old context
and goal are fresh for the newly admitted names, substitute zero for all
of them in the continuation proof $\Gamma,H(\vec b)\vdash A$ and cut
in the proved $H(\vec0)$.  This gives $\Gamma\vdash A$.  The same
argument covers incidental extra parameters in a fractional interval.

These are all successful cases by the root-shape invariant.  The initial
empty-array assumption decodes to the theorem $\top_*$, and ordinary
assumption cut discharges it at the public entry point.
\end{proof}

\subsection{A second, semantic check}

One may instead assign $\bot_*$ to every numeric scalar at a formula
position, leaving its numerical meaning and all array clauses unchanged.
The opening lemma still holds because scalar padding is retained or
replaced by scalar zero, which now has the same falsity reading.  Under
this optional completion, a successful numeric-major call would already
supply native contradiction, and explosion would justify the parent
conclusion irrespective of the false-sentinel match.  Nonempty and empty
array majors have the soundness proof just given.

Thus bottom completion gives a second sufficient semantic argument.
The successful-trace invariant is stronger: numeric majors are actually
unreachable, so changing padding to falsity is unnecessary.  Neither
argument asserts that raw padding is part of the native object language
or that every raw expression has an internally provable equivalence with
its assigned decoder meaning.

\begin{theorem}[Canonical native consequences are preserved]
Every finite native derivation has a certificate in the mandatory-major
presentation.  On canonical native contexts and conclusions the new
checker presents the same consequence relation as $\mathcal A$.
Its ordinary realization is equiconsistent and bi-interpretable with
full first-order iPRA through the inherited whole-domain interpretations.
\end{theorem}
\begin{proof}
First apply the inherited finite repair of excluded computation leaves
using the independently proved nonempty zero initializer.  Take the
canonical backward packet certificate of the resulting derivation.
Keep every actual major and all canonical formulas. Replace each
absent-major marker by the empty array.  In an introduction copy the old
case/progress choice into $u$; in a computation retain the nonempty
initializer.  The added initial truth remains in every retained context.

Reserialize the resulting finite proof in the new context, pruning
any subtree whose goal already closes as a hypothesis, and allocating
names for the remaining proof. This is not an assertion
that every arbitrary historical raw packet stream survives unchanged.
Each mandatory empty-major call is now a hypothesis leaf and consumes
no instruction packet.  All actual majors, computation instances, and
introduction obligations have their original native meanings.  Induction
on the finite native derivation proves successful checking in the
unbounded-resource mathematical reading.  The transformation adds
checker calls at absent-major sites but adds no instruction packets;
it is not claimed to reduce runtime or serialized proof size.

Soundness and this transformation give exactly the native consequence
relation on canonical expressions.  The identity on native syntax
transfers the inherited full numerical and formula interpretation
composites.  The numerical domain remains the whole domain, and the
composite isomorphism has graph $x=y$.  The raw completion is an
implementation semantics for certificate data, not an additional
interpreted object sort or a separate raw-syntax bi-interpretability
claim.
\end{proof}

The new theorem permits a uniform mandatory-major instruction by proving
that its unwanted scalar branch cannot succeed.  The arithmetic basis,
logical reconstruction, and full-iPRA bridge remain inherited.  This is a
structural proof-presentation reduction, not an absolute minimality
theorem or a new arithmetic foundation.  The written trace induction is
not a proof-assistant certificate.
